\chapter{Research Methodology}

\section{Research Workflow}
\TemplateTodo{Explain the research stages from beginning to end. Ensure the workflow aligns with the problem statements and reported results.}

The methodology should be detailed enough for the study to be understood and, where possible, replicated. Describe important decisions, parameters, assumptions, and reasons for method selection.

\section{Research Object and Data}
\TemplateTodo{Describe the research object, data sources, inclusion/exclusion criteria, collection period, data quantity, and ethical or permission aspects if relevant.}

\section{Requirements Analysis}
\TemplateTodo{Describe the hardware, software, and functional requirements. State why each requirement constrains the choices made in this chapter.}

The environment used is specified in Table \ref{tab:requirements}.

% TABLE EXAMPLE 2: some columns fixed, one trailing flexible column.
%
% This is the form you will use most often. A P{0.20\linewidth} column has its
% width locked, while the Y column absorbs whatever page width is left. Put the
% wordiest column in the Y position so the narrow ones are not stretched into
% gappy lines.
%
% Keep the P widths summing to under 1.0\linewidth so the Y column still gets
% room. Below, 0.20 + 0.30 = 0.50, leaving roughly half the page for the notes
% column.
\begin{xltabular}{\linewidth}{|P{0.20\linewidth}|P{0.30\linewidth}|Y|}
\caption{Specification of the computing environment used.}\label{tab:requirements}\\
\hline
Component & Specification & Notes \\
\hline
\endfirsthead
\hline
Component & Specification & Notes \\
\hline
\endhead
\hline
\endlastfoot
Processor & State the model and core count & Explain how it constrains the chosen configuration \\
\hline
Memory & State the RAM capacity & Explain the limit on how much data can be processed at once \\
\hline
Software & State the language, libraries, and versions & Versions are recorded so results can be replicated \\
\end{xltabular}

Two constraints in Table \ref{tab:requirements} bear directly on the research configuration. Explain their consequences here; do not merely repeat the table.

\section{Method or System Design}
\TemplateTodo{Describe the design of the method, model, system, or experimental procedure. Connect each component with the research objectives.}

% \Placeholder marks prose that is not written yet. The box shows up in main.tex
% but disappears in proposal.tex and skripsi-draft.tex, so the manuscript you
% hand to an advisor stays clean. Set \ThesisShowHelpers in
% config/thesis-config.tex to force it on or off.
\Placeholder{Architecture details are still waiting on the first pilot run.}

% \FigurePlaceholder stands in for an image that does not exist yet. Unlike
% deleting the whole figure environment, this keeps the figure number, caption,
% and \label alive, so \ref calls elsewhere do not shift once the real image
% arrives.
\begin{figure}[H]
\centering
\FigurePlaceholder{4cm}{Diagram of the proposed system design}
\caption{Proposed system design}
\label{fig:system-design}
\end{figure}

The system design is shown in Figure~\ref{fig:system-design}. Replace \verb|\FigurePlaceholder| with \verb|\includegraphics| as soon as the image is ready.

\section{Evaluation Scenario}
\TemplateTodo{Explain the experimental design, data split, baselines/comparators, evaluation metrics, hardware/software, and how fair comparison is ensured.}

\section{Data Analysis Technique}
\TemplateTodo{Explain how results will be analyzed: descriptive statistics, metric comparison, error analysis, user validation, or other suitable approaches.}
